\section{Area-aware proof search}
\label{sec:search}
% DRAFT (Claude): awaits the author's rewrite. Matches configs/search.yaml,
% src/markov_lean/search/best_first.py and src/markov_lean/area_prior.py.

\paragraph{Search.}
We use best-first search over goal states in the style of GPT-f and HTPS
\citep{polu2020generative,lample2022hypertree}. Expanding a node asks the
prover, DeepSeek-Prover-V1.5-RL \citep{xin2025deepseekproverv15}, for 8
tactic samples (temperature 1.0, top-$p$ 0.95), runs each in Lean, and pushes
every new goal state with score $\log \pi(a \mid s)$ added to its parent's,
where $\pi$ is a prior over the candidates that produced a new state. With
the baseline prior, $\pi$ is the softmax of the language model's
log-probabilities. The prover is a 4-bit quantization (GGUF, Q4\_K\_M) run by
llama.cpp on an Apple M2 with 16 GB of memory, so absolute solve rates are
below published ones and are not comparable with them. Each theorem gets at
most 12 expansions, a depth of 16 and 180 seconds of wall time, with one
seed. Every condition sees the same tactic samples at each node, so the
comparison between priors is paired.
% TODO(Somesh): add the 64-expansion miniF2F run (configs/search_eval_mac_b64.yaml:
% 64 expansions, depth 32, 1200 s) when it ends.

\paragraph{Area prior.}
The area prior reweights the same candidates by the area each one leads to.
A goal state $s$ is featurized as in Section~\ref{sec:method} and
soft-assigned to its nearest microstates; pushing those weights through the
PCCA+ membership matrix $\chi$ gives a distribution $a(s)$ over the $m$
areas. The target for the next step is
\begin{equation}
  g(s) \propto \bigl(a(s)\,\T_{\mathrm{c}}\bigr) \odot (1 + v)^{-\alpha}
  \odot e^{-\beta\,o},
\end{equation}
where $\T_{\mathrm{c}}$ is the coarse-grained transition matrix, $v$ counts
how often the search has visited each area, and $o$ is the fraction of the
last $w$ expansions spent in each area, applied only when those expansions
closed no goals. The first factor follows the dynamics the model learned,
the second favors areas the search has explored less, and the third pushes
the search out of an area where it has stalled, in the spirit of
metadynamics. A candidate leading to state $s'$ scores $a(s')^{\top} g(s)$;
a candidate that closes the goal scores 1. The scores are normalized to a
distribution $p_{\mathrm{area}}$ and mixed with the language model,
$\pi = (1-\lambda)\,p_{\mathrm{LM}} + \lambda\,p_{\mathrm{area}}$, with
$\lambda = 0.3$, fixed before any search run. The runs here use the first
factor alone ($\alpha = \beta = 0$). The prior is built from the label-free
model of Section~\ref{sec:experiments}, whose large macrostate holds most
steps, so $a(s)$ differs little between candidates: the experiment tests
that model as a search signal, not whether knowing a goal's area helps.

\begin{algorithm}[t]
  \caption{Best-first search with the area prior.}
  \label{alg:search}
  \begin{algorithmic}[1]
    \REQUIRE theorem $s_0$, prover $p_{\mathrm{LM}}$, coarse matrix $\T_{\mathrm{c}}$, mix $\lambda$
    \STATE $Q \gets \{(s_0, 0)\}$; $v \gets \mathbf{0}$
    \WHILE{$Q \neq \emptyset$ and budget remains}
      \STATE pop the highest-scoring $(s, c)$ from $Q$
      \STATE $g \gets$ target for $s$ from $a(s)$, $\T_{\mathrm{c}}$, $v$ and the stall term
      \FORALL{tactic $t$ among 16 samples from $p_{\mathrm{LM}}(\cdot \mid s)$}
        \STATE run $t$ in Lean; skip it on error or a repeated state
        \IF{$t$ closes the goal}
          \STATE \textbf{return} the proof
        \ENDIF
        \STATE $\pi(t) \gets (1-\lambda)\,p_{\mathrm{LM}}(t \mid s) + \lambda\,p_{\mathrm{area}}(t)$, with $p_{\mathrm{area}} \propto a(s')^{\top} g$
        \STATE push $(s', c + \log \pi(t))$ onto $Q$; update $v$
      \ENDFOR
    \ENDWHILE
    \STATE \textbf{return} failure
  \end{algorithmic}
\end{algorithm}

% Shown only once `markov-lean paper` has written figures/search_ablation.pdf,
% which it does when a search run has the area-prior condition (one bar says
% nothing). No synthetic preview stands in for it.
\IfFileExists{figures/search_ablation.pdf}{%
\begin{figure}[t]
  \centering
  \maybefigure{search_ablation}{\columnwidth}
  % TODO(Somesh): rewrite in your words. Drawn by `markov-lean paper`: left, share solved;
  % right, paired change from LM only with the counts solved only by that prior / only by LM.
  \caption{Problems solved under the same budget with the language-model
  prior alone and with the area prior (left), and the paired change from the
  LM prior (right; $+a\,{-b}$: solved only with that prior, only with the LM
  prior). Error bars are 95\% bootstrap intervals over problems. Small local
  budget: \resSearchExpansions{} expansions, \resSearchWallTime{}~s per problem.}
  \label{fig:search}
\end{figure}}{}

\paragraph{Results.}
On the \resSearchProblemsMiniF{} miniF2F test problems the language model
alone solves \resSearchBaseMiniF\% and the area prior \resSearchAreaMiniF\%,
a paired change of \resSearchGainMiniF{} points
[\resSearchGainLoMiniF, \resSearchGainHiMiniF]; random scores at the same
weight change it by \resSearchRandomGainMiniF{} points
(\IfFileExists{figures/search_ablation.pdf}{Figure~\ref{fig:search}, }{}Table~\ref{tab:search}). On the
\resSearchProblemsPutnam{} PutnamBench problems every condition solves the
same \resSearchBasePutnam\%, a floor at this budget that cannot separate the
priors; \resSearchStatementErrorsPutnam{} statements fail to elaborate in
every condition. The area prior from this model does not change search, as
expected from a model that does not separate areas.
% TODO(Somesh): a prior built from judge or premise areas would test the area
% question itself; not run.

\begin{table}[t]
  \caption{Problems solved (\%) under the same budget, \resSearchSeeds{}
  seed. The area prior uses the first factor of $g$ alone at
  $\lambda = 0.3$; random scores keeps the mixing rule and weight with random
  scores in place of the area scores.}
  \label{tab:search}
  \centering\small
  \resultstable{table_search}
\end{table}
